\documentclass[tikz,border=8pt]{standalone}
\usepackage{tikz}
\usetikzlibrary{calc,arrows.meta}
\usepackage{amsmath}
\usepackage{pgfplots}
\pgfplotsset{compat=1.18}
\usepackage{ctex}

\begin{document}
\begin{tikzpicture}

% ===== 左上：无偏性 =====
% 正态曲线（sigma²=0.8），中心 theta=0，峰值约 0.446
% xmax=3.5（不扩展），避免注释延伸进右图空间
\begin{axis}[
  name=ax1,
  width=7.0cm, height=5.8cm,
  xmin=-3.5, xmax=3.5,
  ymin=-0.05, ymax=0.65,
  xlabel={},
  ylabel={密度},
  title={\textbf{无偏性} $E[\hat\theta]=\theta$},
  axis lines=left,
  xtick={-3,-2,-1,0,1,2,3},
  xticklabels={,,,$\theta$,,,},
  xticklabel style={font=\small},
  yticklabel style={font=\scriptsize},
  ytick={0,0.1,0.2,0.3,0.4,0.5},
  grid=major, grid style={gray!20},
  clip=true
]
  % 多次抽样的分布曲线（居中于 theta=0）
  \addplot[blue!70!black, ultra thick, domain=-3.5:3.5, samples=80]
    {exp(-x^2/(2*0.8))/sqrt(2*3.14159*0.8)};
  % 宽的浅色曲线
  \addplot[blue!30, thin, domain=-3.5:3.5, samples=60]
    {0.6*exp(-(x-0.0)^2/(2*1.5))/sqrt(2*3.14159*1.5)};
  % 垂直线标注真值
  \addplot[red!80!black, very thick, dashed] coordinates {(0,0) (0,0.446)};
  % theta 标签：在垂线顶部正上方
  \node[font=\small, red!80!black, above] at (axis cs:0, 0.448) {$\theta$};
\end{axis}

% 注释已整合进子图标题，无需重复

% ===== 右上：一致性 =====
\begin{axis}[
  name=ax2,
  at={(ax1.east)}, anchor=west, xshift=2.0cm,
  width=7.0cm, height=5.8cm,
  xmin=0, xmax=130,
  ymin=0.0, ymax=3.5,
  xlabel={样本量 $n$},
  ylabel={估计量 $\hat\theta_n$},
  title={\textbf{一致性} $\hat\theta_n \xrightarrow{P} \theta$},
  axis lines=left,
  xtick={10,30,50,100},
  xticklabel style={font=\scriptsize},
  yticklabel style={font=\scriptsize},
  ytick={1,2,3},
  yticklabels={$\theta-1$,$\theta$,$\theta+1$},
  grid=major, grid style={gray!20},
  clip=false
]
  % 真值线
  \addplot[red!80!black, very thick, dashed, domain=0:110] {2.0};
  % 真值标签：右侧扩展区
  \node[font=\small, red!80!black, anchor=west] at (axis cs:112, 2.0) {$\theta$};
  % 主估计量曲线（收敛）
  \addplot[blue!70!black, thick] coordinates {
    (5,2.8)(10,1.5)(15,2.6)(20,1.8)(25,2.4)(30,1.9)
    (40,2.3)(50,1.95)(60,2.15)(70,2.05)(80,1.98)(90,2.02)(100,2.0)
  };
  % 第二条曲线（浅色）
  \addplot[blue!30, thin] coordinates {
    (5,1.0)(10,2.9)(15,1.4)(20,2.5)(25,1.7)(30,2.2)
    (40,1.85)(50,2.1)(60,1.95)(70,2.05)(80,2.02)(90,1.99)(100,2.0)
  };
  % n→∞ 注释：下方空白（y in [0.05,0.4]）
  % 此范围内两条曲线均高于 y=1，故 y<0.5 区域为空白
  \node[font=\scriptsize, blue!60!black, align=center]
    at (axis cs:18, 0.35) {$n\to\infty$};
  \draw[->, thick, blue!60!black]
    (axis cs:28, 0.35) -- (axis cs:50, 1.93);
\end{axis}

% ===== 左下：有效性 =====
\begin{axis}[
  name=ax3,
  at={(ax1.south)}, anchor=north, yshift=-1.8cm,
  width=7.0cm, height=5.8cm,
  xmin=-3.5, xmax=4.5,
  ymin=-0.05, ymax=0.95,
  xlabel={估计量取值},
  ylabel={密度},
  title={\textbf{有效性}：方差最小者最优},
  axis lines=left,
  xtick={-3,-2,-1,0,1,2,3},
  xticklabels={,,,$\theta$,,,},
  xticklabel style={font=\small},
  yticklabel style={font=\scriptsize},
  ytick={0,0.2,0.4,0.6,0.8},
  grid=major, grid style={gray!20},
  clip=false
]
  % 高效估计量（方差小，尖峰）
  \addplot[blue!70!black, ultra thick, domain=-3.5:3.5, samples=80]
    {exp(-x^2/(2*0.3))/sqrt(2*3.14159*0.3)};
  % 高效标签：右侧扩展区 x in [1.0, 3.5]，曲线在 x>1 时极小
  \node[font=\small, blue!70!black, anchor=west]
    at (axis cs:1.0, 0.70) {$\hat\theta_1$（高效）};
  % 低效估计量（方差大，宽）
  \addplot[orange!80!black, thick, domain=-3.5:3.5, samples=80]
    {exp(-x^2/(2*1.5))/sqrt(2*3.14159*1.5)};
  % 低效标签：右侧区，避开曲线（x=2.5 时低效曲线约 0.07）
  \node[font=\small, orange!80!black, anchor=west]
    at (axis cs:2.2, 0.17) {$\hat\theta_2$（低效）};
  % 真值线
  \addplot[red!80!black, dashed, very thick] coordinates {(0,0) (0,0.73)};
\end{axis}

% ===== 右下：充分性 =====
\begin{axis}[
  name=ax4,
  at={(ax2.south)}, anchor=north, yshift=-1.8cm,
  width=7.0cm, height=5.8cm,
  xmin=0, xmax=6.0,
  ymin=0, ymax=3.2,
  xlabel={},
  ylabel={},
  title={\textbf{充分性}：统计量无损压缩数据},
  axis lines=none,
  clip=false
]
  % 原始数据框（左侧上方）
  \node[draw=blue!50, thick, rounded corners=4pt, fill=blue!10,
        minimum width=2.0cm, minimum height=0.85cm, font=\small]
    at (axis cs:1.1, 2.5) {原始数据};
  % 数据标签：框下方，不与框重叠
  \node[font=\scriptsize] at (axis cs:1.1, 2.1) {$X_1,\ldots,X_n$};

  % 充分统计量框（右侧上方）
  \node[draw=green!60!black, thick, rounded corners=4pt, fill=green!15,
        minimum width=2.2cm, minimum height=0.85cm, font=\small]
    at (axis cs:4.9, 2.5) {充分统计量};
  % 统计量标签：框下方
  \node[font=\scriptsize] at (axis cs:4.9, 2.1) {$T(X)$};

  % 参数框（正中下方）
  \node[draw=red!60!black, thick, rounded corners=4pt, fill=red!10,
        minimum width=2.0cm, minimum height=0.85cm, font=\small]
    at (axis cs:3.0, 0.8) {参数 $\theta$};

  % 水平箭头：数据框右边 → 统计量框左边
  \draw[-{Stealth[length=6pt]}, thick, blue!70!black]
    (axis cs:2.2, 2.5) -- (axis cs:3.8, 2.5);
  % 箭头上方 T(·) 标签
  \node[font=\scriptsize, blue!60!black]
    at (axis cs:3.0, 2.85) {$T(\cdot)$};

  % 统计量框 → 参数框（右上 → 左下）
  \draw[-{Stealth[length=6pt]}, thick, green!60!black]
    (axis cs:4.5, 2.1) -- (axis cs:3.55, 1.2);

  % 数据框 → 参数框（左上 → 右下）
  \draw[-{Stealth[length=6pt]}, thick, blue!50!black]
    (axis cs:1.5, 2.1) -- (axis cs:2.45, 1.2);

  % 信息等价标签：两箭头中间，不与箭头重叠
  \node[font=\scriptsize, align=center, gray!50!black]
    at (axis cs:3.0, 1.75) {信息等价};

\end{axis}

% 总标题
\node[font=\large\bfseries] at ($(ax1.north)!0.5!(ax2.north)+(0,1.0cm)$)
  {估计量评价四大准则};

\end{tikzpicture}
\end{document}
